\onehalfspacing
\noindent\textbf{\Huge Instructor's signature}

\vspace{3cm} 

\noindent \begin{tabular}{@{}p{0.6\textwidth}@{} p{0.3\textwidth}@{}}
\hrulefill & \hspace{0.3em}\textbf{Date:} \hrulefill \\
\end{tabular}

\begin{flushleft} 
Assoc. Prof. Quan Thanh Tho, Ph.D (Project Instructor)\\
Associate Professor\\
Faculty of Computer Science and Engineering
\end{flushleft}

\newpage

\noindent\textbf{\Huge Declaration of Authenticity}
\vspace{0.5cm}

\noindent My team, which consists of Cao Ngoc Lam and Nguyen Chau Hoang Long, hereby certifies that this specialized project was entirely composed and implemented by us under the guidance and supervision of Assoc. Prof. Quan Thanh Tho at the Faculty of Computer Science and Engineering, Vietnam National University - Ho Chi Minh City University Of Technology.
\vspace{0.5cm}

\noindent Throughout the project, the team ensures that all references and external sources utilized are appropriately cited and thoroughly documented.
\vspace{2cm}

\begin{flushright}
Ho Chi Minh City, December 2025\\
\textbf{The authors},\\
\textit{Cao Ngoc Lam} - \textit{Nguyen Chau Hoang Long}\\
\end{flushright}

\newpage 

\noindent\textbf{\Huge Acknowledgement}
\vspace{1cm}

\noindent First and foremost, we would like to express our sincere gratitude to Assoc. Prof. Quan Thanh Tho for his invaluable guidance, support, and encouragement throughout the development of this project. His expertise and patience have been instrumental in shaping the direction and quality of our work.\\

\noindent We would also like to thank the Faculty of Computer Science and Engineering at Vietnam National University - Ho Chi Minh City University of Technology for providing us with the necessary resources and an inspiring academic environment. Our appreciation extends to all the professors and staff in the faculty who have contributed to our learning and growth during our studies.\\

\noindent Finally, we would like to express our heartfelt thanks to our families and friends for their continuous support, encouragement, and understanding throughout this project. Their guidance and companionship have been a great source of motivation, helping us to complete this work successfully.

\vspace{1cm}

\begin{flushright}
Ho Chi Minh City, December 2025\\
\textbf{The authors},\\
\textit{Cao Ngoc Lam} - \textit{Nguyen Chau Hoang Long}\\
\end{flushright}

\newpage

\noindent\textbf{\Huge Abstract}
\vspace{1cm}

In recent years, programming skills have become increasingly essential for students pursuing careers in technology and software development. However, many learners in Vietnam face challenges in accessing structured platforms that provide comprehensive coding practice, performance evaluation, and competitive benchmarking against peers. Existing platforms available worldwide are not fully tailored to the Vietnamese educational context, often providing contest formats, feedback mechanisms, and progress tracking systems that do not align with local students’ needs, limiting their ability to effectively measure and enhance their programming skills.\\

To address these needs, we developed CodeRank, an online platform designed specifically to support Vietnamese technology students in enhancing their programming capabilities. CodeRank offers a complete environment for creating and participating in coding contests, submitting solutions for automated evaluation, tracking individual progress, and viewing real-time rankings. By leveraging PostgreSQL for reliable data management, the platform ensures consistency, integrity, and efficient handling of multiple simultaneous users, while the responsive frontend delivers a smooth and intuitive user experience.\\

The system is developed using a modular approach, breaking the project into smaller components to facilitate incremental progress and maintainable code. Technologies such as ReactJS for the frontend and NestJS for backend services were chosen for their flexibility, scalability, and active community support. Advanced features, including performance analytics and dynamic ranking dashboards, provide students with actionable insights to guide their learning and foster a competitive yet educational environment. Additionally, an AI-powered chatbot has been integrated into the system, serving as an interactive interview coach for algorithmic coding challenges. Leveraging Google Gemini as the underlying large language model and LangChain as the orchestration framework, the chatbot simulates real interview scenarios, maintains contextual memory through a hybrid short-term and long-term mechanism, provides step-by-step guidance, and delivers constructive feedback tailored to each student’s responses. This feature enhances the learning experience by enabling learners to practice effectively, receive hints, and refine their problem-solving skills in a realistic, guided environment.\\

Overall, CodeRank presents a structured, interactive, and educational platform that empowers Vietnamese technology students to develop programming skills, engage in coding competitions, and monitor their growth, contributing to improved learning outcomes and readiness for professional challenges.
